\section{Efficient Endpoint Calling with API Speculative Execution}
\label{sec:endpoint_speculative}

The previous sections introduced optimizations for open-weight models, which require modifications to the model inference process. Additionally, \lang works with API-access-only models, such as OpenAI's GPT-4. However, for these models, we can only call a black-box API endpoint.

This section introduces a new optimization for black-box API models that accelerates execution and reduces the API cost of multi-call \lang programs using speculative execution. For example, a program may ask the model to generate a description of a character with a multi-call pattern: \texttt{s += context + "name:" + gen("name", stop="\textbackslash n") + "job:" + gen("job", stop="\textbackslash n")}. Naively, the two \texttt{gen} primitives correspond to two API calls, meaning that the user needs to pay for the input token fee on the \texttt{context} twice.
In \lang, we can enable speculative execution on the first call and let it continue the generation of a few more tokens by ignoring the stop condition.
The interpreter keeps the additional generation outputs and matches and reuses them with later primitives.
In certain cases, with careful prompt engineering, the model can correctly match the template with high accuracy, saving us the latency and input costs of one API call.


